\section{Two numbers under every smile}
\label{sec:fwdtwo}

Part~I built three families of smile models, and all three lived in the
same coordinates.  A strike $K$ entered as its log-moneyness
$k=\log(K/F)$; a dollar price entered divided by $DF$
(\cref{sec:intro}).  Two numbers sat under every
computation: the \emph{forward} $F$, the price agreed today for
delivery of the underlying at expiry, and the \emph{discount factor}
$D$, the value today of one dollar paid at expiry.  Part~I treated
both as given.  This chapter asks where they come from.

The honest answer is that, for a single stock or an ETF, neither is
quoted.  There is no screen showing ``the forward of SPY for December
2026''.  There is a spot price, and there is the option chain itself:
for each strike, a call and a put.  The forward and the discount must
be \emph{estimated} from those option prices --- which makes them
fitted quantities, with error bars, exactly like the smile parameters
they will later support.  That change of status is the subject of the
chapter, and it opens Part~II's theme: before a model can be fitted,
the observation itself has to be built, and every step of that
construction is either measured from the data or chosen and disclosed.

The estimate rests on one identity.  Fix an expiry and a strike $K$,
and write $\mathsf{C}(K)$ and $\mathsf{P}(K)$ for the dollar mid
prices of the call and the put there --- sans-serif letters, here and
for the rest of the book, mark a raw market quote, as opposed to the
normalized model prices $c$ and $P$ of Chapter~2.  Assume for now that both options are \emph{European} --- exercisable
only at expiry, in the standard vocabulary; \cref{sec:fwdline} meets
the American reality.  Consider the portfolio that is long the call
and short the put, and follow it to expiry.  If the underlying finishes at $S_T$ above $K$, the call pays
$S_T-K$ and the put pays nothing; if it finishes below, the call pays
nothing and the short put costs $K-S_T$.  In both cases the portfolio
delivers
\[
(S_T-K)^+-(K-S_T)^+=S_T-K ,
\]
the payoff of a forward contract struck at $K$.  What is that worth
today?  Price the two pieces separately.  The constant $-K$ is easy: a
sure payment of $K$ at expiry is worth $DK$ now, so the piece
contributes $-DK$.  For the $S_T$ piece, use the definition of the
forward price: a forward contract struck at $F$ costs \emph{nothing}
to enter --- that is what makes $F$ the forward price --- and it pays
$S_T-F$.  Prices of payoffs add, so the value of receiving $S_T$
minus the value of receiving the constant $F$ equals the value of
receiving $S_T-F$, which is zero.  Receiving $S_T$ is therefore worth
exactly what receiving $F$ is worth: $DF$.  The portfolio's price today must therefore be
\begin{equation}
\Pi(K)\;\equiv\;\mathsf{C}(K)-\mathsf{P}(K)\;=\;D\,(F-K).
\label{eq:fwdparity}
\end{equation}
This is \emph{put--call parity} \citep{Stoll1969,Hull2018}.  Chapter~2
already carried its normalized form --- the rightmost identity of
eq.~\ref{eq:calldef}, $c(k)-P(k)=1-e^k$; multiplying it by $DF$ and
writing $K=Fe^k$ recovers \eqref{eq:fwdparity} exactly.  Note what the
derivation did \emph{not} use: no model of how the underlying moves,
no volatility, no distribution.  Parity is a validity statement in the
sense of the book's thesis --- any European call and put prices that
violate it admit an arbitrage, whatever the model.

Now read \eqref{eq:fwdparity} the other way.  The right-hand side is a
straight line in $K$, with slope $-D$ and root $F$.  The left-hand
side is observable: one number per strike, computed from quotes.  So
the chain itself draws a line, and the two numbers Part~I assumed are
that line's two coefficients.  Estimating $(F,D)$ means fitting a
straight line --- nothing more exotic than that.

Before any general machinery, solve the smallest possible case by
hand.  Suppose an expiry half a year away quotes two usable strikes,
and the portfolio prices come out as
\[
\Pi(90)=11.27,\qquad \Pi(110)=-8.33
\]
(constructed numbers, chosen to come out exactly; appendix
\ref{app:fwdprotocol}).  Two points determine the line.  The slope is
\[
\frac{\Pi(110)-\Pi(90)}{110-90}
=\frac{-8.33-11.27}{20}
=\frac{-19.60}{20}
=-0.98 ,
\]
so the discount is $D=0.98$.  The intercept --- the line's value at
$K=0$ --- is $\Pi(90)+D\cdot90=11.27+88.20=99.47$, and by
\eqref{eq:fwdparity} the intercept equals $DF$, so
\[
F=\frac{99.47}{0.98}=101.50 .
\]
Finally, the discount converts to an interest rate.  Writing $t$ for
the \emph{calendar} year fraction to expiry ($t=0.5$ here) and
$D=e^{-rt}$,
\[
r=-\frac{\log D}{t}=-\frac{\log 0.98}{0.5}\approx 4.04\% .
\]
Two portfolio prices --- four option quotes --- have produced a
forward, a discount, and an interest rate, out of prices alone.  This is worth pausing on: the option
market carries its own funding information, and the fitter never needs
an external rate feed to normalize a smile.

One clock deserves a sentence before we go on.  The book has used
$\tau$, variance time, since Chapter~2, and this chapter introduces
$t$, the calendar year fraction.  They are different objects with
different jobs: variance accrues on the market's clock (Chapter~8
makes that precise), but discounting and \emph{carry} --- the growth
rate that links spot to forward, given its full content in
\cref{sec:fwddivs} --- accrue on the calendar: a dollar earns
interest on weekends too.  In everything below, $D=e^{-rt}$ and the
carry exponents use $t$; total variance $w=\sigma^2\tau$ keeps
$\tau$.  In this chapter's data the two clocks
happen to agree numerically, so nothing subtle hides in the figures;
the distinction starts to matter in Chapter~8.

We now know what the two numbers are and that the chain determines
them.  The next question is what happens with a real chain: many
strikes, quoted with noise, where the line must be fitted rather than
solved.
